\section{Conclusion}
\label{sec:conclusion}

We revisited evidence selection in multimodal medical RAG and argued that the
prevailing recipe---score candidates for relevance and take the top $K$ in a
single pass---is fundamentally passive: it conflates relevance with
answer-utility, offers no recourse when the retrieved pool is inadequate, and,
on visually homogeneous medical images, systematically admits high-confidence
but clinically wrong evidence. We recast selection as an \emph{agentic search}
process. A vision-language agent judges image and text candidates cross-modally
and either accepts them for generation or rewrites the query to re-retrieve,
trained with an \emph{answer-utility reward}---the lift in a frozen generator's
probability of the correct answer---which needs only multiple-choice supervision
and no passage-level labels. A density-aware router allocates the retrieval
budget on every round and can suppress a modality whose signal is noise, and we
evaluate under a corrected end-to-end image--text protocol that feeds the
retrieved images to the generator directly.

Our framework has limits. The reward is a per-query answer-probability lift, so
it presumes access to gold answers at training time and is most natural for
multiple-choice supervision; extending answer-utility rewards to open-ended
medical generation, and reducing the cost of iterative re-retrieval, are natural
next steps. More broadly, we hope answer-utility--driven, cross-modal agentic
retrieval proves useful beyond endoscopy---wherever visual similarity is a poor
proxy for the evidence a model actually needs.
